L'outil auquel nous avons réfléchi permet donc aux utilisateurs de se préparer en amont de leurs recherches avec les outils de consultation. Cet outil leur permet de comprendre et de voir la structure des bases de données du point de vue du \gls{td}. De cette façon, il leur est plus rapide de trouver les notices des documents contenant les informations qu'ils désirent. Ils peuvent réunir des corpus de document qu'ils désirent consulter en sachant exactement quelles informations il y trouvera. De cette façon, les utilisateurs de cet outil sont moins dépendants à la fois des connaissances des documentalistes et des outils vieillissants de l'INAthèque ne correspondant plus aux modes de consultations actuels. 